\section{Scope, relation to the repository, and principal results}
\label{sec:scope}

The canonical manuscript in \cite{repo} already develops a substantial
periodic Hurwitz--Stieltjes convolution calculus. In particular, its
Fourier-normalized family on $\R/\mathbb Z$ leads to shifted correlations,
finite parts, coincident limits, contact terms, and primitive ladders.
Repeating those formulas under new notation would not advance that program.
Here the integration variable is instead the imaginary part of the
\emph{second} Hurwitz argument:
\begin{equation}\label{eq:vertical-scope}
 \frac1{2\pi}\int_{-\infty}^{\infty}
       f(a+\ii y)g(b-\ii y)\dd y,
 \qquad \Rea a,\Rea b>0.
\end{equation}
This is neither an integral in the spectral variable $s=\sigma+\ii y$ nor
a periodic correlation on $[0,1]$. It requires a different subtraction,
a different transform, and different convergence hypotheses. The auxiliary
Hurwitz products over $[0,1]$ studied in \cite{shpot2016,shpot-paris} also
have a different integration geometry and use a different auxiliary
function. Our starting transform identities are classical; the contribution
of this note is the explicit, jointly proved closure calculus and its
specializations. We do not assert priority for every displayed identity.

Set
\begin{equation}\label{eq:Z-definition-intro}
 Z_s(z)=\zeta(s,z)-\frac{z^{1-s}}{s-1},
 \qquad
 G_m(z)=\gamma_m(z)+\frac{\log^{m+1}z}{m+1}.
\end{equation}
The value at $s=1$ is removable, and $G_0(z)=\log z-\psi(z)$.
The main equal-scale theorem evaluates \eqref{eq:vertical-scope} for
$f=Z_s$, $g=Z_t$ by a one-variable function $K(s+t-1,a+b)$.
Its all-order Taylor coefficients then evaluate every pairing of $G_m$
and $G_n$. Finite Bernoulli subtractions give the same conclusion for
negative spectral orders. Rational vertical dilations introduce only a
finite set of shifted Hurwitz arguments, rather than an irreducible double
zeta object.

Three consequences indicate the scale and normalization of the results.
Define the analytic Stirling remainder in the right half-plane by
\begin{equation}\label{eq:Omega-intro}
 \Omega(z)=\log\Gamma(z)-(z-\tfrac12)\log z+z-\tfrac12\log(2\pi).
\end{equation}
Then
\begin{align}
 \frac1{2\pi}\int_{\R}
   |G_0(\tfrac12+\ii y)|^2\dd y
   &=\log(2\pi)-\frac32,\label{eq:headline-digamma}\\
 \frac1{2\pi}\int_{\R}
   |\Omega(\tfrac12+\ii y)|^2\dd y
   &=\zeta'(-1)+\frac{\zeta(3)}{2\pi^2}+\frac19,
       \label{eq:headline-stirling}\\
 \frac1{2\pi}\int_{\R}
    G_0(\tfrac12+\ii y)G_0(\tfrac12-2\ii y)\dd y
   &=\log\Gamma(\tfrac14)-\tfrac12\log\pi-\tfrac14\log2
       \notag\\[-1mm]
   &\hspace{9mm}+\tfrac34\log3-\tfrac34-\tfrac\pi8.
       \label{eq:headline-quarter}
\end{align}
Their right-hand sides are approximately $0.33787706640934548$,
$0.0065868815274467235$, and $0.22363092143574916$, respectively.
These numerical values are diagnostics; the proofs appear below.

\subsection{What is and is not resolved}
The questions answered here are precise: whether spatially compensated
vertical Stieltjes pairings close at arbitrary indices; whether negative
spectral orders can be handled without assigning finite parts to the
vertical integral; and whether rationally related slopes preserve finite
Hurwitz closure. The answer to each is affirmative, with explicit formulas
and proofs. These are newly developed continuations in this package, not a
claim that a named conjecture in an unread incoming archive has been solved.
The repository's $S_6$ and revised $S_8$ problems are not settled here.

The audit was targeted. The canonical README, the periodic convolution
foundation, a terminal rational-reduction section, the incoming README,
and incoming metadata were inspected. Not every recent ZIP was decoded or
audited. Thus absence of a formula from these inspected sections is not a
proof of absence from the entire corpus. The package includes an integration
map and a source inventory so that an independent maintainer can perform
that final overlap check.

\subsection{Conventions that affect the answer}
Throughout, $\log z$ is the principal logarithm on $\Rea z>0$.
The function $\log\Gamma(z)$ means its analytic logarithm, real on the
positive real axis; it does \emph{not} mean the pointwise principal value
of $\log(\Gamma(z))$. A prime on $\zeta$ or $\Phi$ denotes differentiation
in the spectral order, unless a different variable is explicitly indicated.
We use
\begin{equation}\label{eq:laurent-convention}
 \zeta(1+u,z)=\frac1u+
      \sum_{m\ge0}\frac{(-1)^m}{m!}\gamma_m(z)u^m.
\end{equation}
All the vertical integrals asserted as theorems are ordinary absolutely
convergent Lebesgue integrals. The words ``constant term'' and ``removable
value'' used later refer to a finite \emph{spectral expression}, not to a
regularization of a divergent integral in $y$.
